For the source comparison, use the intervals~\eqref{eq:E-window}. Fix $b>0$ large enough for the degree and the finite tests, put $l=t-h_t(b)$, $k=l-v_t$ and $r=a_t(\cutoffsize_t)^b$. The interval from $t$ to~$l$ separates the reading points from the last fluctuation scale; the interval from $l$ to~$k$ permits the corrected sources to increase in size while their remainders decrease. With $L=\ell_{J_t}^{\rm st}$, these intervals satisfy
\begin{equation}\label{eq:E-window-geometry}
 L^{h_t(b)}\ge(\cutoffsize_t)^{b+2},\qquad
 L^{v_t}\ge B_t^{10},\qquad
 \widehat\beta_i\in[\widehat\beta_t/2,\widehat\beta_t]\quad(k\le i\le t)
\end{equation}
on sufficiently late intervals. The last assertion follows from $h_t(b)+v_t=o(t)$ at fixed stage and the entry choices. All rounded row points lie strictly inside the level-$k$ cell: their distance from the origin is $O(r)$, while $a_l/r\ge(\cutoffsize_t)^2$. Rounding changes their rescaled coordinates by $O((\cutoffsize_t)^{-b})$.

Fix the arity bound $n\ge2$ and the compact separated position set specified above. Put $\mathcal M=\cutoffsize_t$. Let $q_{\rm der}$ bound the required derivative orders, and let $c,\delta_{\rm G}>0$ be the localization and healthy-chart exponents in the estimates below. Define
\begin{equation}\label{eq:E-source-initial}
 \mathcal E_t^{\rm G}=\mathcal M^{-b}
 +C_j\mathcal M^{q_{\rm der}}
   (\widehat\beta_t^{-\delta_{\rm G}}+e^{-cB_t}).
\end{equation}

The polynomial comparison supplies the initial small error for the subsequent source iteration.
\begin{lemma}[Middle-level source comparison]\label{lem:E-source-polynomial}
For each fixed finite label family, sufficiently late intervals admit common real lifts, fixing the unit and bounded in $N,p$ for that interval. Uniformly over the prescribed compact position set, all corrected nonempty modes at~$l$ have total regular and normalized forced bound
\begin{equation*}
 \mathcal C_j^{h_t(b)+2}\mathcal E_t^{\rm G}.
\end{equation*}
The assertion holds on a full cutoff tail, uniformly in~$p$.
\end{lemma}
The proof gives the Wick correction and the inverse-separation gain.
\begin{proof}

Use Definitions~\ref{def:D04}, \ref{def:D07}, \ref{def:D11}, Lemmas~\ref{lem:C1}--\ref{lem:C2}, \ref{lem:D-connected}, \ref{lem:E-marked-step}, \ref{lem:E-replication} and the covariance estimates of Lemmas~\ref{lem:B2} and~\ref{lem:B-windows}. The sources have at most $n$ marks in each cell; the number of cells is unrestricted. Use the coefficientwise rooted clause of Lemma~\ref{lem:D-connected} in their weighted nilpotent algebra. Only point separation is needed for the covariance comparison; coincident coordinates of distinct points are allowed.


On a fixed healthy chart, subtract the radial Taylor polynomial of each local insertion. The coefficient norm of the remainder is bounded by its next ordered derivative times the corresponding power of the coordinate radius. All derivatives needed here are among the finitely many stage orders. Express each polynomial in curvature and differentiated-curvature variables, carrying along their powers of~$r$. For a polynomial~$\pi$, the finite Gaussian contraction operator is
\begin{equation*}
 \mathsf H_C\pi=\sum_{j\ge0}\frac1{2^jj!}(\Delta_C)^j\pi,
 \qquad \mathsf H_C^{-1}\pi=\mathsf H_{-C}\pi.
\end{equation*}
The sums are finite on the chosen polynomial space, since a contraction removes two factors. Use the inverse operator for contractions internal to a single insertion. This is the Wick-ordering correction; contractions between distinct insertions remain. The leading coefficient change $T_m$ of Lemma~\ref{lem:E-marked-step} and this finite contraction operator, in the jet normalization of Definition~\ref{def:D07}, are invertible and preserve the unit and real form. The error in replacing the selected coefficient map by this leading map has norm less than $1/4$ after its inverse, by the marked estimate~\eqref{eq:E-atom-estimates} and the entry condition. Solving this finite linear system gives one lift for each free basis polynomial, common to every row. Transport only between the fixed levels $r(d)$ and~$t$ bounds these reference coefficients independently of~$N,p$ for that interval. All matrices are real on the real form, and the unit equation has the exact solution~$\mathbf u_d$.

The covariance calculation involves finitely many derivatives of $\Gamma(x)=(4\pi^2|x|^2)^{-1}$. For each pair of coloured plaquette stencils the rescaled Fourier chart $K(\xi)$ is supported in $|\xi|\le2$. Its leading symbol is the Maxwell curvature symbol $H(\xi)$. For every fixed required order~$q_{\rm der}$, the quadratic-form calculation of Lemma~\ref{lem:B2}, differentiated in frequency, gives
\begin{equation*}
 |\partial^jK(\xi)|\le C|\xi|^{-j},\qquad
 |\partial^j(K-H)(\xi)|\le C|\xi|^{1-j}
 \quad(0<|\xi|\le1,\ j\le q_{\rm der}).
\end{equation*}
These bounds follow by differentiating the inverse quadratic symbol: its curl factors remove the inverse-square singularity; replacing each lattice difference by its linear term leaves one extra power of~$|\xi|$. The regulator and its derivatives obey the same fixed bounds; different colours have zero cross covariance. To convert the second estimate to position space, decompose a cutoff of $K-H$ into annuli $|\xi|\asymp2^{-v}$. A derivative of order~$s$ in position inserts $s$ powers of~$\xi$. After $q_{\rm der}>5+s$ integrations by parts on each annulus its transform at~$Rx$, with $x$ in a compact separated set, is bounded by
\begin{equation*}
 C2^{-v(5+s)}\min\{1,(R2^{-v})^{-q_{\rm der}}\}.
\end{equation*}
Splitting the sum at $2^v\asymp R$ gives $CR^{-5-s}$. The smooth part away from zero decreases faster than this. The leading homogeneous symbol is the corresponding curvature derivative of~$\Gamma$ away from the origin, with the Fourier normalization of Part~I. Including the curvature scaling powers therefore gives relative error $O(R^{-1})$.

For the finite jet samples, terms of the matching homogeneity give this leading derivative exactly. Every other term has at least one extra inverse power of~$R$; summing their finitely many absolute coefficients preserves the $O(R^{-1})$ bound. Here $R=(\cutoffsize_t)^b$, and rounding in the rescaled variables gives an additional $O(R^{-1})$. Localizing the sector covariance introduces $C(\cutoffsize_t)^{q_{\rm der}} e^{-cB_t}$ by the inverse-covariance bound in Lemma~\ref{lem:B-windows}, with smaller $c$, since $\widehat\beta_i\ge\widehat\beta_t/2$ throughout the interval. Its nonlinear healthy-chart error is $C(\cutoffsize_t)^{q_{\rm der}}\widehat\beta_t^{-\delta_{\rm G}}$ for some fixed $\delta_{\rm G}>0$. The derivative orders and powers~$q_{\rm der}$ are fixed by the degree and the finite row family; their stage constants are fixed before entry. Since $\log\cutoffsize_t=o(\log\widehat\beta_t)$ and $B_t\asymp(\log\widehat\beta_t)^2$, both errors are smaller than any prescribed fixed negative power of~$\cutoffsize_t$ after entry. Choose $b$ so that the inverse-separation error is smaller than the finite power required by the source radius below. The chart and inspected supports lie inside the nonwrapping comparison interval in~\eqref{eq:E-window-geometry}, so these local estimates hold on every period, with the same constants.


The finite-polynomial estimate combining these errors is as follows. Let $U$ bound the Gaussian comparison polynomials, $T$ the one-factor error, and $Z$ the two-factor error in their coefficient norm. For the connected two-point term, subtract the product of the one-point terms. Expanding that product gives the bound
\begin{equation}\label{eq:E-Gaussian-polynomial}
 T+Z+(U+T)T+TU.
\end{equation}
For example, if $a_i$ and $a_i^0$ are the two one-point terms, then
$a_1a_2-a_1^0a_2^0=(a_1-a_1^0)a_2+a_1^0(a_2-a_2^0)$, whose norm is at most $(U+T)T+UT$. The terms $T,Z$ are the errors in the remaining one- and two-mark polynomial terms. Gaussian contraction expands into pairings, so the bound for any fixed polynomial is a finite polynomial in the bounds for its contractions. After $b$ has been fixed sufficiently large, the preceding covariance and Taylor estimates make~\eqref{eq:E-Gaussian-polynomial} smaller than every source bound required at~$l$.

For more slots, let $\mathcal G'_J$ be the Gaussian polynomial moment before replacing its covariance by the Maxwell covariance. On the compact separated position set, all these moments are bounded by a fixed~$U\ge1$, and their discrepancies are at most $\varepsilon$, a constant times the covariance and Taylor error above. For a partition with $q$ blocks, telescoping its product bounds the discrepancy by $q\varepsilon(U+\varepsilon)^{q-1}$. The partition formula for $\mathsf C^{\rm Max}$ therefore bounds the sum of all nonempty corrected coefficients by
\begin{equation*}
 C_n\varepsilon(1+U+\varepsilon)^{n-1}.
\end{equation*}
There are finitely many subsets and partitions for fixed~$n$. Their constants are included in~$\mathcal C_j$. Subtracting $\mathsf C^{\rm Max}$ cancels every leading free coefficient in the logarithm; for a pair this is precisely the mean-product subtraction in~\eqref{eq:E-Gaussian-polynomial}. A support carries a mode only when it contains all its anchors. The support-independent scalar correction is retained separately per cell. Thus no local term is assigned to a support lacking one of its marks.

For terms meeting a flagged region, keep that region's complete source support and its original eligibility test. At a step $i+1\to i$ in stage~$j$, let $v_0$ be the integrated majorant of the distinguished marked factor and $n_0$ its covering number, as in Lemma~\ref{lem:D-connected}. Its rooted forced bound is $v_0e^{\eta n_0}e^{-r_j\beta'_u t_{i+1}^2}$ before division by the forced source allowance. Half this suppression pays the derivative and support factors of the finite marks; the other half bounds their normalized forced remainder by the error above after entry. The finite polynomial contractions, Taylor subtractions and connected attachments cost at most $\mathcal C_j$ at each of the $h_t(b)$ steps. Their product is the factor in the statement. At each step the source norm is submultiplicative and the rooted sum is anchored at a mark, so the bound is independent of the number of replicated cells.
\end{proof}

The next estimate increases the source radius while retaining the full regular and forced representation.
\begin{lemma}[Corrected Gaussian sources]\label{lem:E-Gaussian-sources}
Uniformly over the prescribed compact position set, arbitrarily late intervals admit common real lifts of the free basis, bounded in $N,p$ for each fixed interval and fixing the unit, such that the flowed corrected source $e^{-\mathsf C^{\rm Max}(z)}\prod_h(1+z_hF_h^{\rm row})$ has a representation~\eqref{eq:E-source-form} at~$k$ with radius
\begin{equation*}
 \alpha_t=\exp(\log\cutoffsize_t/(8n)).
\end{equation*}
At this radius its regular and forced source sums are each at most one, with scalar source bound at most one per cell. The assertion is eventual in~$N$, uniformly in~$p$.
\end{lemma}
Its proof iterates the middle-level error over the second part of the interval.
\begin{proof}

Use Definitions~\ref{def:D04}, \ref{def:D05a}, Lemmas~\ref{lem:D-connected}, \ref{lem:E-marked-step}, \ref{lem:E-replication} and~\ref{lem:E-source-polynomial}. Work in the submultiplicative weighted mode norm of Lemma~\ref{lem:E-replication}, with the same support weights and forced eligibility tests.

We now bound the loss over all steps; their number grows with~$t$. Use~\eqref{eq:E-source-initial}. Increasing the fixed power $q_{\rm der}$ includes the finite Taylor and rounding bounds above. The source radius contributes at most $\alpha_t^n=\mathcal M^{1/8}$ to any mode in one cell. Enlarge $\mathcal C_j$ to include all per-step finite polynomial and support constants. From~\eqref{eq:E-log-cost} and the two ceilings in~\eqref{eq:E-window}, on sufficiently late intervals with fixed~$b$,
\begin{equation}\label{eq:E-source-loss}
 \mathcal C_j^{h_t(b)+v_t+2}
 \le\mathcal M^{(b+3)/j^2}B_t^{10/j^2}
 \le\mathcal M^{1/64}.
\end{equation}
Indeed $\log B_t=2\log\log\widehat\beta_t+o(1)$ and $\log\mathcal M=(\log\log\widehat\beta_t)^2$, so $B_t=\mathcal M^{o(1)}$. The constant factors from the ceilings fit in the extra $\mathcal M^{1/j^2}$; the final inequality holds at all sufficiently large stages for each fixed~$b$.

We use the following one-step source estimate. For the fixed lifts of Lemma~\ref{lem:E-source-polynomial}, let $S_i$ majorize the corrected singleton terms, and $P_i$ the corrected terms with at least two marks, in the weighted regular norm at radius~$\alpha_t$, including their scalar terms per cell. Use the same quantities for the forced norm after division by its positive source allowance, and take the larger values. For $k\le i<l$, the connected product expansion gives
\begin{align}\label{eq:E-source-step}
 S_i&\le K_j^{\rm src} S_{i+1}+\mathcal M^{1/2}S_{i+1}^2,\notag\\
 P_i&\le K_j^{\rm src} P_{i+1}
          +\mathcal M^{1/2}(S_{i+1}+P_{i+1})^2.
\end{align}
Here $K_j^{\rm src}\ge1$ is fixed at the stage, and $\mathcal C_j\ge K_j^{\rm src}+2$. To obtain these bounds, separate connected terms linear in the incoming corrected sources from those with at least two source factors. The former use the one-mark integration bound. In the latter, the product rule leaves a product of two source norms; summing the additional attachments uses the rooted clause of Lemma~\ref{lem:D-connected}. Its finite derivative and localization losses are bounded by $\mathcal M^{1/2}$ after entry. Every higher product is absorbed in the same quadratic bound while the source norm is at most $\mathcal M^{-1/2}$. For the forced terms the retained eligibility event supplies the forced suppression before its allowance is divided out. The entry requirement in item~\ref{it:E-entries} includes this normalized forced bound. Disjoint cell modes use the multiplicative weights of Lemma~\ref{lem:E-replication}; their connecting supports keep the same estimate. No factor counts the number of cells.

Set $\mathfrak a_i=S_i+P_i$. Equations~\eqref{eq:E-source-step} imply $\mathfrak a_i\le K_j^{\rm src}\mathfrak a_{i+1}+2\mathcal M^{1/2}\mathfrak a_{i+1}^2$. The first part of the interval gives $\mathfrak a_l\le\alpha_t^n\mathcal C_j^{h_t(b)+2}\mathcal E_t^{\rm G}$. By~\eqref{eq:E-source-loss}, for $b\ge1$ and sufficiently late entry,
\begin{equation}\label{eq:E-source-iteration}
 \alpha_t^n\mathcal C_j^{h_t(b)+v_t+2}\mathcal E_t^{\rm G}
 \le\mathcal M^{-3/4},\qquad
 \mathfrak a_i\le\mathcal C_j^{l-i}\mathfrak a_l\le\mathcal M^{-3/4}<\mathcal M^{-1/2}.
\end{equation}
The second inequality follows by backward induction: below $\mathcal M^{-1/2}$ the quadratic term in the recurrence is at most $2\mathfrak a_{i+1}$. For the first, the covariance contribution is at most $\mathcal M^{1/8+1/64-b}$, whose exponent is at most $-55/64$. The other contributions tend to zero faster than every fixed negative power of~$\mathcal M$, since $\log\mathcal M=o(\log\widehat\beta_t)$ and $B_t\asymp(\log\widehat\beta_t)^2$. This proves compatibility of the initial-error condition with the whole growing interval, rather than treating its constants as a finite product independent of~$t$.

The extra length $v_t$ also provides the support margin: a level-$l$ neighborhood of radius $B_t^2$ has level-$k$ radius at most $B_t^2/L^{v_t}\le B_t^{-8}$. It therefore fits inside the protected cell around its rounded reading points. The successive kernel neighborhoods add a bounded geometric sum in those units; their fixed margin was included in the choice of the comparison cell. Thus the inspection and attachment tests used in the iteration stay inside the stated interval. The unit has no remainder by Lemma~\ref{lem:E-marked-step}.

The support margin in~\eqref{eq:E-window-geometry} makes every replicated inspected stencil fit its own cell. Terms meeting flagged regions keep their original forced eligibility and the suppression used in Lemma~\ref{lem:D-connected}; their source amplitude replaces only the marked letter's amplitude. Thus the local source sums at the enlarged radius are at most one in both regular and forced norms. The constant source contribution is at most one per cell. These are estimates on the complete source representation, not on its Gaussian polynomial alone. They prove the lemma on arbitrarily late intervals, with all thresholds independent of~$p$.
\end{proof}

The remaining normalization uses the complete positive unmarked density. In particular it cannot discard the forced sum or replace it by its healthy part.
\begin{lemma}[Normalized source integral]\label{lem:E-source-integral}
For the corrected sources of Lemma~\ref{lem:E-Gaussian-sources}, put $u_t=\log\cutoffsize_t/(16n)$. Eventually on its intervals, uniformly in~$p$, for every nonempty mode~$I$,
\begin{equation*}
 \left|\left\langle\prod_{C\in\mathcal C^{\rm tile}_k}(\mathfrak c_I)_C\right\rangle_{N,p}\right|
 \le e^{|\mathcal C^{\rm tile}_k|q_t^{\rm cmp}}
       \alpha_t^{-|I||\mathcal C^{\rm tile}_k|},
 \qquad q_t^{\rm cmp}\le u_t.
\end{equation*}
\end{lemma}
\begin{proof}
Use Definition~\ref{def:D06} and positivity of $\density P$, which follows from the density identity of Theorem~\ref{prop:RG}, the positive Wilson start and the probability kernels of Lemma~\ref{lem:B1}, and Lemma~\ref{lem:E-Gaussian-sources}. Let $P$ be the unmarked representation at the width level~$k$, $V=|\Lambda_{n_{p,k}}|=|\mathcal C^{\rm tile}_k|$, $g=\beta(P)^{-1/2}$, $H_{\rm path}=A_{J_k}^{\rm path}$, and let $\delta,\Lambda^{\rm loop}$ bound its regular and loop terms. The admissibility inequalities give $0<g\le t_k$, $\delta\le1$, $\Lambda^{\rm loop}\le\Lambda_k^{\rm loop}$ and $\beta(P)\le4\widehat\beta_k$.

Choose constants $s_H\in(0,1]$, $b_H\ge1$ such that small Haar balls have measure at least $s_Hr^{\dimG}$ and the clipped one-link integrals are at most $b_H(g/\sigma_*)^{\dimG}$, where $\sigma_*$ is the smaller of one and the linear clip coefficient. These follow from the nonsingular exponential chart and compactness of the group, with constants chosen before all levels. Write
\begin{equation*}
 \mathfrak h=\log b_H-\dimG\log\sigma_*,\qquad
 \mathcal P(P,X,g)=2E_{\rm ex}+6+
 \beta(P)E_{\rm base}+X+3\mathfrak h-4\log s_H.
\end{equation*}
Here $E_{\rm ex}$ bounds the loop and regular exponent per site, including the exponent discrepancy $\delta$ in the clipped comparison, and on links of size at most $g/4$,
$\beta(P)E_{\rm base}=3(c_{{\rm dir},J_k}+H_{\rm path}^2)$. For any integrable complex~$F$ with
$z|F(U)|\le e^{\Phi_P(U)+VX}$, $z\ge0$, the ratio estimate is
\begin{equation}\label{eq:E-positive-ratio}
 \frac{z|\int F\,dU|}{\int\density P\,dU}
 \le e^{V\mathcal P(P,X,g)-4\dimG\log g}.
\end{equation}
For the denominator, fix an axial spanning tree with $V-1$ links. Gauge invariance removes those links, leaving $3V+1$ free links. Restrict each free link to the ball of chord $g/4$. All plaquettes are then healthy, so the forced covering sum equals one. The lower action bound and the Haar-ball bound give
\begin{equation*}
 \int\density P\,dU\ge
 e^{v_{{\rm vac}}V-V(E_{\rm ex}+\beta(P)E_{\rm base})}
 (s_Hg^{\dimG})^{3V+1}.
\end{equation*}
For the numerator, the clipped plaquette inequality is
\begin{equation*}
 \beta(P)E_{\rm plaq}(U)\ge
 \frac{\chord(U_\square)}{g/\sigma_*}-1.
\end{equation*}
It follows in the quadratic branch from $(a/g-1)^2\ge0$ and $\sigma_*\le\min(1,c^{\rm dir})$, and in the linear branch from $\sigma_*\le c^{\rm lin}$ and $g$ below the threshold. Summing over the six plaquettes per site retains the product of $e^{-\chord(U_\square)/(g/\sigma_*)}$ in the upper integrand.

Keep the $3V-3$ plaquette factors selected as follows and discard the others, whose values lie in $[0,1]$. The axial tree returns each site towards zero in its highest nonzero coordinate. Order free links by that coordinate, treating wrap links first in each coordinate, except for the four links at zero, and then transverse links in decreasing height. Associate to each treated link its plaquette towards the tree. That plaquette contains the link once; every other remaining selected plaquette is independent of that link in this order. Haar invariance therefore replaces its factor by the one-link integral. The four untreated links account for the difference $(3V+1)-4=3V-3$. All discarded plaquette factors lie in $[0,1]$, including at the periodic boundary. Thus
\begin{equation*}
 z|\int F\,dU|\le
 e^{v_{{\rm vac}}V+V(E_{\rm ex}+6+X)}
 \bigl(b_H(g/\sigma_*)^{\dimG}\bigr)^{3V-3}.
\end{equation*}
Divide these inequalities. The powers of $g$ differ by $-4\dimG$. Using $3V-3\le3V$, $3V+1\le4V$, $b_H\ge1$ and $s_H\le1$ gives~\eqref{eq:E-positive-ratio}. This proves the estimate for the complete positive unmarked denominator without removing its forced terms outside the restricted lower-bound set.

One may use the clipped comparison exponent in this calculation even if it is not the exponent used for the source bound. Their difference has absolute value at most $V\delta$. Multiplying the upper source bound by $e^{V\delta}$ and bounding the two positive partition integrals in the opposite direction by the same factor accounts for this change. These two factors are included in $2E_{\rm ex}$ because each occurrence of $E_{\rm ex}$ includes $\delta$. Thus~\eqref{eq:E-positive-ratio} applies to the actual source integral on the original exact history.

The enlarged-radius bounds of Lemma~\ref{lem:E-Gaussian-sources} give $X\le1+n$: one is the scalar bound, and each of at most $n$ marks contributes at most one. The local coefficient estimate before integration includes $z=\alpha_t^{|I|V}$. The entry inequalities imply
\begin{equation*}
 E_{\rm ex}\le8H_{\rm path}+2,\qquad
 -4\dimG\log g/V\le1.
\end{equation*}
For the second, $-4\dimG\log g=2\dimG\log\beta(P)$ is $O(\log\widehat\beta_k)$, whereas $V\ge D_k=e^{s_k}$. The first follows from
$\widehat\beta_kt_k^3\le1$, $\widehat\beta_kt_k^4\le1$ and $\Lambda_k^{\rm loop}t_k^2\le1$ on sufficiently late levels. Hence choose a stage constant
\begin{equation}\label{eq:E-comparison-cost}
 Q_j^{\rm cmp}\ge2(8A_j^{\rm path}+2)+6+3(c_{{\rm dir},j}+(A_j^{\rm path})^2)+(1+n)
 +3\mathfrak h_j-4\log s_H+1
\end{equation}
for each arity admitted at stage~$j$. It bounds the exponent per cell in~\eqref{eq:E-positive-ratio}. All its quantities are fixed before the entry to stage~$j$. Item~\ref{it:E-entries} in Lemma~\ref{lem:E1.1} requires $Q_j^{\rm cmp}\le\log\cutoffsize_t/(16n)$ on the chosen intervals. This proves $q_t^{\rm cmp}\le u_t$ and the asserted coefficient estimate. In particular the volume appears only through the number of cells, so the cutoff threshold is independent of~$p$.
\end{proof}
